% !TEX program = xelatex
% 编译方式：xelatex main.tex （连续编译两次以生成目录；建议用 latexmk -xelatex）
% 依赖宏包：ctex / xeCJK / tcolorbox / enumitem / booktabs / array / colortbl / xcolor / fancyhdr / hyperref

\documentclass[UTF8,a4paper,12pt,fontset=windows]{ctexart}

% ---------- 页面与行距 ----------
\usepackage[a4paper,left=2.6cm,right=2.6cm,top=3.0cm,bottom=3.2cm]{geometry}
\renewcommand{\baselinestretch}{1.32}
\setlength{\parskip}{0.2em}

% ---------- 中文字体 ----------
\setCJKfamilyfont{yahei}{Microsoft YaHei}
\newcommand{\yh}{\CJKfamily{yahei}}

% ---------- 常用宏包 ----------
\usepackage{amsmath,amssymb}
\usepackage{enumitem}
\usepackage{booktabs}
\usepackage{array}
\usepackage{longtable}
\setlength{\tabcolsep}{4pt}
\emergencystretch=3em
\raggedbottom
\usepackage{xcolor}
\usepackage{tcolorbox}
\tcbuselibrary{breakable,skins}
\usepackage{fancyhdr}
\usepackage[colorlinks=true,linkcolor=teal,urlcolor=teal]{hyperref}

% ---------- 配色 ----------
\definecolor{ink}{RGB}{40,44,52}
\definecolor{teal}{RGB}{15,84,96}
\definecolor{clay}{RGB}{186,98,74}
\definecolor{mist}{RGB}{243,246,244}

% ---------- 章节标题 ----------
\ctexset{
  part = {
    format = {\centering\yh\huge\bfseries\color{teal}},
    aftertitle = {\par\vspace{4pt}{\color{clay}\rule{0.4\linewidth}{1.4pt}}\par},
    beforeskip = {3em},
    afterskip = {2em},
  },
  section = {
    format = {\yh\Large\bfseries\color{teal}},
    beforeskip = {2.4em},
    afterskip = {0.8em},
  },
  subsection = {
    format = {\yh\normalsize\bfseries\color{ink}},
    beforeskip = {1.6em},
    afterskip = {0.5em},
  },
  subsubsection = {
    format = {\yh\normalsize\bfseries\color{ink}},
    beforeskip = {1.2em},
    afterskip = {0.4em},
  },
}

% ---------- 自定义组件 ----------
% 核心洞察框
\newtcolorbox{insight}{
  breakable, enhanced,
  colback=mist, colframe=teal, boxrule=0.8pt,
  arc=1.5mm, left=2.2mm, right=2.2mm, top=1.6mm, bottom=1.6mm,
  before skip=6pt, after skip=8pt
}
% 提醒 / 边界框
\newtcolorbox{notehot}{
  breakable, enhanced,
  colback=white, colframe=clay, boxrule=0.8pt,
  arc=1.5mm, left=2.2mm, right=2.2mm, top=1.6mm, bottom=1.6mm,
  before skip=6pt, after skip=8pt
}
% 清单环境（□ 复选框）
\newenvironment{checklist}
  {\begin{itemize}[leftmargin=2em,itemsep=1.5pt,topsep=3pt,label={$\square$}]}
  {\end{itemize}}
% 术语强调
\newcommand{\term}[1]{\textbf{\color{teal}#1}}
% 核心一句
\newcommand{\key}[1]{\par\vspace{2pt}\noindent\textcolor{teal}{\textbf{核心一句：}}#1\par\vspace{4pt}}

% ---------- 页眉页脚 ----------
\pagestyle{fancy}
\fancyhf{}
\setlength{\headheight}{16pt}
\fancyhead[L]{\small\yh\color{ink}最舒适的认知：ADHD 自救讲义}
\fancyhead[R]{\small\yh\color{ink}\leftmark}
\fancyfoot[C]{\small\color{ink}\thepage}
\renewcommand{\headrulewidth}{0.3pt}

\hypersetup{pdftitle={最舒适的认知：ADHD 自救讲义},pdfauthor={自我交付版}}

\begin{document}

% ---------- 封面 ----------
\begin{titlepage}
  \centering
  \vspace*{2.2cm}
  {\yh\Huge\bfseries\color{teal} 最舒适的认知\par}
  \vspace{0.5em}
  {\yh\LARGE\bfseries\color{ink} ADHD 自救讲义\par}
  \vspace{1.8em}
  {\large\color{ink} 从学历上层到学历中下层\par}
  {\large\color{ink} 习惯 · 成因 · 方法论 · 思考\par}
  \vspace{2.8cm}
  {\color{ink} 思想内核提炼自 B 站「阿卡迪萨 · ADHD 野路子」系列 24 期视频\par}
  {\color{ink} 在其自成一派的猜想之上，延伸为一份完整的认知与自救讲义\par}
  \vfill
  {\large 2026 年 \quad 自我交付版\par}
  \vspace{1.5cm}
  {\color{clay} 写给每一个觉得自己「明明可以，却总做不到」的人\par}
\end{titlepage}

% ---------- 目录 ----------
\pagestyle{empty}
\pagenumbering{roman}
\tableofcontents
\clearpage
\pagestyle{fancy}
\pagenumbering{arabic}

% ---------- 正文 ----------
\section*{写在前面：这份讲义是什么，不是什么}

先花三分钟，把话说清楚。

这份讲义是写给谁的？写给成人 ADHD、自认 ADHD、以及正怀疑自己可能是 ADHD 的成年人的。你不需要诊断书，也不需要用「没确诊」来否定自己——这本讲义讨论的是「这套大脑怎么运转」，诊断与否，不影响它对你有多少用。

这份讲义从哪来？它的思想内核，提炼自 B 站 UP 主「阿卡迪萨」（ADHD 野路子）的 24 期视频。一个清华数学博士，在 B 站搞笑的过程中发现自己疑似 ADHD，然后把自己的大脑当成一台机器来拆、来猜、来调试，留下了一整套自成一派的猜想。这份讲义的最大特点，正是继承了那个态度：\term{不讲医学权威，只讲自己大脑的实际运行日志}。

那这份讲义不是什么？它不是医学诊断，也不是吃药指南。它不承诺治愈。它只做一件事：\term{让你第一次觉得，自己没坏}。

\key{最舒适的认识，一句话：你不需要把自己修成正常人。你需要的是读懂自己的参数，然后顺着参数生活。}

所谓「舒适」，不是躺平。舒适是停止和重力对抗——你不再天天质问自己「为什么别人能做到我却不能」，而是开始研究「我这套系统到底吃什么燃料、走什么路线、在什么条件下运转得最好」，然后照着这个去设计生活。

讲义的路线图，一共四部分：

\begin{checklist}
  \item \textbf{底层认知}：ADHD 到底是一套什么样的算法，它真正的麻烦在哪；
  \item \textbf{理论延伸}：同一个算法，为什么在学历上层和学历中下层会走出两种人生；
  \item \textbf{自救方法论}：先让认知舒服，再给一套可操作的方法库，最后从你的位置出发；
  \item \textbf{更深一层的思考}：一个更舒服的看待世界、看待自己的角度，和一页纸的每日清单。
\end{checklist}

文末还有两份附录：一份给「怀疑自己是 ADHD」的人的自检线索，一份本讲义的黑话速查表，方便你随时回来翻。

最后，一个必要的声明，也是原视频反复强调的态度：\term{以下是猜测，个体差异极大}。所有方法都是「试试看」，不是「必须做」。好用就用，不好用就扔。这份讲义里的每一个办法，原样照搬大概率会失灵，按自己的参数改装之后才可能灵。

如果你已经被生活严重困住，该寻求专业帮助就去找，这不丢人，也和读这份讲义不冲突。但在这之前、在这之后，都请先记住这句话——你不是懒，你只是还没被读懂。

\part{底层认知：ADHD 是一套算法，不是一处故障}

\section{ADHD 是一套算法，不是一处故障}

先聊命名。\term{「注意力缺陷多动障碍」}这个译名，把所有人都带偏了——它让你以为 ADHD 就是「注意力不够、太多动」。于是 ADHD 本人也天天自责：「我要是能再专注一点就好了。」但真正的图景，和「专注力不够」几乎没有关系。

原视频给了一个更好的框架，我们称它为\term{算法观}（BV13HoGYiEgF）：

把大脑理解成一套算法。普通大脑的成本结构是：深度运算（单点纵深、线性推进）便宜，主动注意力便宜；而 ADHD 大脑恰好反过来——\term{广度运算便宜，深度运算贵；被动注意力便宜，主动注意力贵}。

这不是「坏了」，是「成本结构不同」。就像 GPU 和 CPU：GPU 跑并行便宜，跑单线程反而贵；CPU 正好相反。你不能指着 GPU 说「你怎么这么不会跑单线程」，它是另一种东西。

顺着这个框架，很多现象第一次能解释清楚了：

\begin{itemize}
  \item \term{为什么能刷六个小时视频，却不能看十分钟书}——刷视频用的是便宜的被动注意力，看书用的是昂贵的主动注意力。
  \item \term{为什么上课听不懂就彻底崩了}——课堂要求连续线性跟随，那是深度运算，是 ADHD 最贵的项目。
  \item \term{为什么「重要」的事情推不动}——「重要」是理性的语言，而你的启动器只听多巴胺的语言。
\end{itemize}

\subsection{椭圆竞争与超专注}

注意力空间里，无意识的念头（原视频叫「椭圆」）一直在互相竞争。普通人的主动注意力圈大，能把冒出来的念头按住；ADHD 的主动圈小，念头经常赢，这就是走神。

但当某一个念头强到占满几乎整个注意力空间，就会出现\term{超专注}——连续几小时沉浸，外界全部屏蔽。注意，这里的关键不是「能不能专注」，而是\term{电源开关不在你手上}：你没法在「该专注的时候」打开它，只能等它自己来。

超专注常见于游戏，但它可以被引导。原视频分享过一段真实经历：大三上学期，他能一天学十三个小时，是因为他当时真的相信「把定义和定理一个个弄清楚，一步步推，就能考进前几名」——\term{他看到了真凉}（实利）。「想考好成绩」这个念头足够强，于是在注意力竞争里赢了，连续数月霸占整个空间。

\begin{insight}
ADHD 自救的核心，从来不是「把主动注意力练大」（那是逼 GPU 去跑 CPU 的活，又贵又伤），而是\term{让被动注意力更多地落在正事上}——把那个免费、便宜、量大管饱的通道，引导到对的地方。原视频叫它\term{以广促深}。
\end{insight}

\subsection{核心特征与伴随特征：别把所有毛病都往 ADHD 上贴}

网上关于 ADHD 的描述，九成在讲坏处，而且讲得五花八门：记性差、粗心、忘时间、情绪化、拖延……看起来每一条都像「你」。原视频把这里叫\term{归并效应}（BV1zU9LBVEnm）：一个标签被广泛使用之后，那些标签并不必然导致的东西，也会被贴到这个标签上。

但真实情况是：\term{ADHD 的核心特征只有一样——执行功能缺陷}（启动困难、维持线性集中困难）。其余全是伴随特征，个体差异极大，甚至互相矛盾：有人记性差得要命，有人对感兴趣领域的知识记得极牢；有人情绪化，有人反而理性得过分。

所以「我是不是真的 ADHD」这个问题，本身可能问错了。更值得问的是：\term{这套算法，在我的生活里造成了多少实际损害？}按 DSM 的标准，确诊 = 特征 + 功能损害，两者缺一不可。同一个人，在清华当数学博士时每天写数学写得好好的，没法确诊；在博士期间连最简单的任务都启动不了、每天五点就想睡觉时，一测一个准。\term{ADHD 是一个「环境交互」的标签，不是一个贴在脑门上的本质属性}。

\begin{notehot}
这一段不是让你别去了解自己，恰恰相反——是让你别被「网上的 ADHD 画像」吓到或框死。你可以有 ADHD 的算法，同时记性极好；你可以没确诊，但这套方法对你也照样有效。标签只是地图，不是领地。
\end{notehot}

\key{你不是不能专注，你是「主动专注」这项服务对你来说太贵了。与其硬充，不如学会用便宜的通道。}

\section{真正的麻烦：无聊、低唤醒、动机、主动切换}

上一章说完了「算法观」，这一章回答一个更实际的问题：ADHD 真正卡住你的，到底是哪几件事？

答案是四件——\term{无聊、低唤醒、动机不足、主动任务切换成本过高}。专注力不够反而是最表面的那个。逐一说。

\subsection{无聊是一种认知负担，不是闲得慌}

普通人对「分心」的默认态度是：分心是敌人，要消灭它。但原视频给出的观察正好相反：\term{对 ADHD 来说，无聊才是大问题，分心不是}（BV1q9PmzxEt2）。

硬坐、硬逼自己专注，对神经连接的负担，比「带着分心做事」大得多。你逼自己盯着一页纸，脑子里念头一个接一个往外冒，每压一个念头都在耗电——而那个念头本来可能是回血用的。

所以 ADHD 的首要任务不是「别分心」，是\term{先别无聊}。原视频甚至建议：先别管定向，优先保证「一天能做得久、做得开心」。无聊解决了，一天干十几个小时也不累；定向的调控，是后面的、漫长的、慢慢练的事。

\subsection{低唤醒：人醒着，脑子部分睡着了}

ADHD 经常经历一种很诡异的疲劳：明明没干什么重活，却觉得浑身没力气；午休睡了两个小时，起来反而更累。这不是矫情，也不是懒，原视频把它解释为\term{低唤醒状态}（BV16mXnBxEE1）——有研究显示，ADHD 成人在清醒状态下，大脑会出现类似睡眠的慢波活动。也就是说：\term{你看着是醒着的，但有一部分脑子已经睡着了}。

这个现象带来两个重要推论：

\begin{itemize}
  \item \term{休息对 ADHD 是「唤醒恢复」，不是「静止恢复」}。被动躺平、午睡、硬睡，对普通人是充电，对 ADHD 可能是放电——没有刺激，大脑进入更昏沉的状态。这也是为什么刷几分钟视频、看到个感兴趣的内容，你反而清醒了。
  \item \term{低唤醒时，执行功能下线，但自然语言和审查能力还在}。你没法线性地看一段有难度的内容，但可以「随便瞄一下」，然后冒出几个问题去问 AI——这是 ADHD 专属的低配模式，照样能干活。
\end{itemize}

针对唤醒度，原视频给了一套三档管理法：

\begin{checklist}
  \item \textbf{高唤醒}：做线性任务，趁热打铁；
  \item \textbf{中唤醒}：随手翻找，看有没有什么激发兴趣；
  \item \textbf{低唤醒}：躺下，允许大脑自发放出想法，用手机或 AI 里的内容重新唤醒，而不是硬撑着执行。
\end{checklist}

\key{你觉得自己「没体力」的时候，往往不是真的没体力，是唤醒度太低——缺的不是休息，是刺激。}

\subsection{意志耗能：抑制一个念头，可能比做那件事更费电}

原视频有个非常反直觉的观察（BV1DQAkzbEuK）：\term{压住一个念头去做别的事，消耗的能量可能比做事本身还大}。

他举的例子是拍视频：脑子里一个视频「读条读满了」，很想录，但按老规矩「早上录完，下午做数学」，于是得把录视频的冲动压一个下午——结果数学也没做成，还多睡了两小时。后来改成「随时写、随时录」，抑制的耗能消失，数学和视频两不误。

把这个放大到生活：取个快递本身三分钟体力，但你拖着不去，是因为出门可能撞见同事、而工作还没做完、怕被评判——\term{压住「怕被评判」的念头，比取快递费电得多}。执行功能在这里扮演的角色，与其说是「干活」，不如说是「维持 + 意志」：维持住目标，压住其他念头。

\begin{insight}
所以 ADHD 的调控方向，不是「提高意志力」，而是\term{降低需要意志的场合}。怎么降？原视频给了一个关键原则：\term{降低权重，而不是排序}。别给自己排一个 ABCD 的执行顺序表——排序本身就是高耗能的抑制过程。把一堆东西在内心合并成一件大事（后面「恰当场」会细讲），让大脑觉得「今天我就做这一件事，其他都不重要」，意志耗能立刻降下来。
\end{insight}

\subsection{主动切换贵，被动切换便宜}

这是整个方法库的地基（BV11ucZzxEkx）：

\term{主动任务切换}（我正做着 A，强迫自己去做 B）＝ 抑制当前冲动，贵得离谱；
\term{被动任务切换}（脑子里的念头自己冒出来，跟着它走）＝ 念头自发涌现，便宜。

所以最优策略不是「按计划切换」，而是\term{把一堆事合并成一条宽的抽象轨道，让念头在轨道里自己跳}。外部看你在频繁切换，内部看你一直在一件事上。这一条的实操化，就是后面方法论里的「恰当场」。

\subsection{假性疲劳：缺的是定向能量，不是追随性能量}

最后一种假性疲劳（BV1MWwJzxEh5）：你累的时候，缺的往往是\term{定向能量}（「我接下来必须写那篇东西」），而不是\term{追随性能量}（「随便看看有什么有意思的」）。

很多人拖着不做正事时的循环是这样的：五点多回宿舍，一想「晚上要看一篇文章」，立刻觉得「好累，先睡一会」，于是从五点多睡到七点、八点、九点……一晚上就这么过去了。\term{「必须要做 X」这个念头本身，就在前额叶持续耗电}，你休息得越久，越起不来。

破解思路是\term{动态调节自由度}：想法上高自由度（「今天不一定要写这个，随便偏出去写点别的也算成果」），操作上低自由度（把要写的文件放在桌面最显眼处，记录一直开着）。这样既取消了「必须」的警报，又保证「啪的一下」的想法能最快落到正事上。

\key{ADHD 的问题，大部分不是意志力问题，是能量管理问题。先把无聊和焦虑解决了，再谈定向。}

\section{动机悖论与冲动理性：你的奖励系统长在别处}

上一章讲了能量，这一章讲动机——为什么「重要」的事推不动，而「没意义」的事反而写得飞快。

\subsection{动机悖论：压力越大，越启动不了}

原视频讲过一个让很多人拍大腿的例子（BV1iKPrzNEjN）：公司让你写一本书，告诉你写好了能卖钱——你一直拖。但如果换个念头：「人类要完蛋了，这本书写不写无所谓，随便写写图个好玩」——你反而很快就写完了，质量还不错。

从逻辑上看这是个悖论：你都觉得人类要完蛋了，还写什么书？但从神经机制上看，一点都不悖论：

\term{「能卖钱、要交差、写不好会被评判」这些社会规范压力，会推高你的皮质醇。皮质醇一高，启动就锁死了。而「人类要完蛋」这个念头，功能上消除了社会规范压力，皮质醇降下来，你反而能动了。}

原视频为此给了一个极重要的心智工具：\term{一元论视角}。不存在一个坐在大脑里发号施令的「统一的自我」；信念不是指令，而是网络状态。一个念头（「人类要完蛋」）的价值不在于它是不是真理，而在于它产生了什么效果——系统只问效用，不问真理。只要它好用，就把它留着，不必要求自己逻辑自洽。

\begin{insight}
这对 ADHD 自救的意义非常大：\term{别执着于「我应该是自洽的、认真的、负责的」}。那些宏大叙事——「我要做一个靠谱的人」「这份工作必须一丝不苟地完成」——很可能正是锁死你启动的东西。该自我欺骗的时候就自我欺骗，只要它真的降低焦虑、释放执行力，从神经学的角度看，它就是真实。
\end{insight}

\subsection{冲动理性 vs 时间理性}

普通人的「理性」，是一种\term{时间理性}（BV1apQ2ByEN3）：假设干 A 还是干 B，区别只有那点切换成本，所以按部就班永远最优。

但 ADHD 的「理性」不是这样运行的。原视频用打游戏举例子：你砍一个血很厚的 BOSS，砍几刀就觉得无聊想走——时间理性会说「走了还得重新加载，浪费几十秒，不划算」；但\term{冲动理性}会说：「无聊了就走，因为现在有个更有价值的冲动冒出来了，冲动伴随新想法，走是对的，BOSS 以后再来打。」

对 ADHD 来说，\term{价值权重不在「这个程序花多少时间」，而在「当下的神经兴奋度，以及兴奋指向哪里」}。同样一份小组报告，冲动来的时候三十分钟写完；被硬安排着写，一天都写不完，还要边写边想「这破报告对我有啥用，写不好还要被嘲笑」。

所以：

\begin{checklist}
  \item \term{冲动不是敌人，是信号}。第一次冲动大多没用，但它是种子——二次冲动、三次冲动会长成深刻的问题。随手记录下来，别让它溜走（工具上 Obsidian 非常合适）。
  \item \term{把冲动引导成网，而不是捋成线}。你不可能把冲动排成一条直线，但可以引导成一张大致有用的网——在大方向内跳来跳去，而不是被任务死死锁死。
  \item \term{考试能专注的真相}：考试时你走不了，环境强制压制了其他冲动。但这个办法没法日常化——现代人有电脑有手机，把自己关进小黑屋也有一百种办法分心。
\end{checklist}

\subsection{兴趣引导：为什么「无敌」之后反而无聊}

单机游戏开挂变成无敌，很多人反而觉得没意思了——这几乎是所有人的体验，但原视频指出它和 ADHD 有深刻的联系（BV1Y2wizyEua）：\term{大脑在扫描「有没有值得做的事」，扫描结果是没有，于是陷入无聊}。游戏无敌后是这样；ADHD 面对「写了也没用」的任务时，也是这样。

区别在于：普通人能说服自己「这东西能卖钱、能交差」，硬着头皮驱动自己；ADHD 对这些外部驱动力不那么感冒，大脑的预测结果就是「这东西写了也没啥用」，于是干脆无法启动。

\begin{insight}
所以 ADHD 的正确姿势不是「逼自己启动一个写了也没用的任务」，而是\term{让自己尽可能待在一个可以跳跃的激发场里}——你定向驱动的能量非常有限，但你长期经营好一个场域，让有意义的问题在这个场域里冒出来的概率足够高，然后顺着它做。这叫\term{以广促深}，是贯穿这份讲义的总心法。
\end{insight}

\subsection{隐形回血：管理分心，而非消除分心}

最后补一个被低估的机制（BV16hxrzYEpT）：\term{全神贯注的效率，会被事后累抵消}。

让 ADHD 全神贯注地敲代码，任务期间确实更快，但敲完可能要躺三个小时——把效率全赔回去。反过来，允许一些「有价值的创新性分心」：敲代码时突然冒出一个想法，「这个定理如果和那个东西混合会怎么样」，花一分钟记下来或问一下 AI——这不只是回血，还是好想法的种子，甚至会让「启动这件事」的门槛降低（因为你知道中途可以合法地跑神）。

所以 ADHD 的目标不是\term{消除分心}，是\term{管理分心}——把分心的范围控制在有价值区间内：

\begin{checklist}
  \item \textbf{无新内容的冲动}（刷手机、开无关网页、想断连）：用身份认同压下去——「我是一个能按顺序行动的人」；
  \item \textbf{有点意思的想法}（冒出一个新的连接、新的猜想）：记录下来或问 AI，那是回血，也是种子。
\end{checklist}

\key{对 ADHD 谈「效率」，常常是复收益：没有快速反馈时，效率指标里全是「做了也没用」的评判。真正的奖励不是钱，是随机冒出来的好想法——那才是你的充电桩。}

\part{理论延伸：同一个算法，两种剧本}

\section{学历上层：被聪明和脚手架掩盖的 ADHD}

先讲一个反直觉的事实：\term{学历上层反而是 ADHD 最隐蔽的地方}。

\subsection{画像}

一路名校，硕士博士，体面的职业。小时候「看着做的事挺多、成绩还行」，从没被任何人怀疑过。可能直到二十多岁甚至三十多岁，在某次崩溃、某段关系破裂，或某个视频下面，才第一次认真怀疑：我是不是 ADHD？

为什么这么晚才发现？因为他们的算法缺陷，被两样东西盖住了：\term{聪明}，和\term{脚手架}。

\subsection{他们是怎么「混」过来的}

\begin{itemize}
  \item \term{智力补丁}：低要求阶段，靠聪明直接覆盖执行缺陷。小学初中内容简单，被动注意力够用，不用主动注意力也能学个七七八八。
  \item \term{肾上腺素冲刺}：靠死线启动。D 前三天是天才，D 前三个月是废物——这不是夸张，这是常态。紧迫感是唯一有效的启动器，于是每件事都被他们亲手拖到变成紧急事。
  \item \term{超专注红利}：偶尔对某件事上头，一次超专注顶普通人几个月的线性积累。以广促深在长周期里确实成立，所以他们能「平时不学，考前封神」。
  \item \term{环境脚手架}：好学校、好家庭、好导师，处处有人替他们把结构架好。课表有人排，作业有人收，死线有人催。他们的执行功能缺憾，被环境无声地覆盖了。
\end{itemize}

\subsection{习惯观察}

学历上层的 ADHD，日常往往是这样一副矛盾的样子：

\begin{checklist}
  \item \textbf{死线前爆发}；能拖则拖；对「无趣但有长期回报」的任务完全无感；
  \item \textbf{生活层面的功能崩塌}：堆积的快递、没拆的包装、未回的微信、未缴的账单、乱到没法收拾的房间。工作体面，生活失能；
  \item \textbf{自我药物化}：咖啡因、酒精、短视频、游戏、报复性熬夜——靠刺激续命；
  \item \textbf{完美主义面具}：越聪明越不敢交「不完美」的东西，于是更拖。别人催他，他说「还没准备好」，其实是在等那个永远等不到的「准备好了」。
\end{checklist}

\subsection{成因：环境把执行缺陷盖住了}

这里要说透一个扎心的事实：\term{他们不是没问题，是问题被环境买断了}。学校的课表、导师的催逼、家庭的后勤，都是外部脚手架——是环境在替他们的前额叶上班。这解释了为什么很多高功能 ADHD 一脱离校园（进入自由度极高的科研、自由职业，或独居生活）就立刻崩掉：脚手架撤了，缺陷第一次裸奔。

\begin{insight}
学历上层的核心矛盾：\term{能力天花板很高，燃料系统却在慢性透支}。他们是靠肾上腺素而不是靠系统运转的——每一次冲刺都是透支，多巴胺阈值越抬越高，于是越来越「撑不住」。很多人的「中年危机式的觉醒」，本质就是这个引擎烧到冒烟了。
\end{insight}

\subsection{代价：燃尽、冒充者、与越成功越空心}

\begin{itemize}
  \item \term{燃尽循环}：冲刺—透支—恢复—再冲刺，周而复始，直到某次彻底转不动。高功能 ADHD 最常见的就医触发点，不是「做不好」，而是「突然做不动了」。
  \item \term{冒充者综合征}：总觉得自己是「靠运气和最后一晚混过来的」，随时可能被拆穿。这个想法不完全是错的——他们确实在靠最后一晚——但代价是不敢停、不敢露怯、不敢求助。
  \item \term{越成功越空心}：成绩、职位带来的多巴胺在边际递减。小学考第一还有奖励，读完博士发完顶刊，奖励系统彻底失灵——「然后呢？」「这有啥用？」这正是第三章讲的动机悖论在顶端的显影。
  \item \term{恐惧失败导致不敢开始}：因为太聪明，太知道「这件事如果认真做能做成什么样」，反而被那个标准吓住，于是宁可拖着。
\end{itemize}

\subsection{一句话总结这一章}

学历上层不是「没事」，是\term{一直有事，只是还没崩}。他们的自救起点，不是「学会更努力」，而是承认自己需要脚手架——哪怕外表看起来一切正常。

\section{学历中下层：被系统和标签漏掉的 ADHD}

如果说学历上层是「被聪明掩盖」，那么学历中下层就是\term{被系统漏掉}——而且在被漏掉的过程中，还被打上了一层最毒的标签。

\subsection{画像}

也许从未确诊，甚至根本不知道 ADHD 是什么；早早离开学校，或勉强读完技校、大专；做着需要动手、需要即时反馈、不需要长时间线性规划的工作；也可能长期在家，状态低迷，被家人说「懒」、被自己说「废」。

他们中的很多人，在互联网上第一次看到 ADHD 的描述时，是哭着看完的——不是因为确诊了，而是因为终于有一个词，能解释自己这么多年的「不对劲」。

\subsection{他们经历了什么}

理解学历中下层的 ADHD，必须从学校说起。因为\term{学校可能是全世界对 ADHD 最不友好的环境}：

\begin{itemize}
  \item \term{教育系统默认深度优先算法}：一节课四十五分钟、一章接一章、知识线性积累、一次没跟上就永远跟不上。这恰好是 ADHD 算法最贵的运行模式，和它是天生的仇家。
  \item \term{早期失败，归因到态度而不是大脑}：「不认真」「懒」「笨」「态度有问题」。归因到态度，比归因到大脑更伤人——因为它意味着「你本来可以，是你不愿意」，是把「不能」翻译成了「不愿」，再翻译成「人品问题」。
  \item \term{自我概念崩坏}：在最重要的成长期，被反复告知「你不行」，于是「我做不到」成了默认信念，哪怕他其实很聪明。
  \item \term{回避 → 更多失败}：逃避所有可能暴露「我不行」的场合——不交作业、不去考试、不竞选、不发言——于是错过更多机会，失败进一步坐实。
  \item \term{家庭资源不足}：没有诊断的机会、没有好老师、没有「再试一次」的试错资本。学历的差距会滚雪球：一步落下，后面每一步都在为前面买单。
\end{itemize}

\subsection{习惯观察}

学历中下层的 ADHD，在外人眼里常常是这样：

\begin{checklist}
  \item \textbf{三分钟热度}：兴趣一来就上头，几天几夜不睡觉；兴趣一退就弃坑，换下一个。频繁换工作、频繁开始又放弃；
  \item \textbf{用即时刺激填满生活}：游戏、短视频、烟酒——在旁人看来是堕落，其实是一个燃料不足的大脑在拼命找刺激自救；
  \item \textbf{隐藏技能}：极强的动手能力、直觉、口才、销售天赋、危机处理能力——这些在学校里不被计分，在社会上也不被当成「学历」，但它们真实存在；
  \item \textbf{「地下」的成功模式}：很多人最终在不需要学历的赛道里反而活得不错——手艺、销售、自由职业、创业。这正是广度算法的用武之地。
\end{checklist}

\subsection{成因：不是能力不够，是失败的积累压垮了执行功能}

要特别强调一件事：\term{学历中下层的 ADHD，核心问题从来不是「能力」}。同一颗大脑，放进一个有人推着走的环境，能读出来；放进一个没人管、没人信的环境，可能连毕业都难。

他们真正缺的是两样：\term{一次成功的经验}（「我能做到」的证据），和\term{一个不评判他的环境}。前者是燃料，后者是脚手架。这两样都没有，执行功能就会在「反正我也做不到」的预测里彻底躺平。

\begin{insight}
学历中下层最缺的不是方法，是\term{「我能做到」的证据}。方法可以学，但若心里那个「我不行」的信念不破，什么方法都会在第一次失败时被当成「你看，果然不行」。所以他们的自救顺序，和学历上层正好相反：学历上层先补系统，学历中下层先修自我概念。
\end{insight}

\subsection{一句话总结这一章}

学历中下层的核心矛盾：\term{不是做不到，是相信自己做不到}。他们最需要的是被给一次「第二次机会」——不管是自己给的，还是世界给的。

而这一章写在这里，就是那份第二次机会的第一步：\term{你的学历低，不代表你不行；它只代表，你一直没有脚手架。}

\section{学历的两面：它度量的是什么}

把前两章放在一起看，会得出一个让人舒服、也让人清醒的判断：

\key{学历度量的大多是「你拿到了多少外部脚手架」，而不是「你有多大的内在能力」。}

同样是广度算法的大脑，放进有老师推、有家长托、有资源垫的环境里，能走出高学历；放进没人管、没人信的环境里，可能连高中都读不完。这不是能力之差，是环境之差。

\subsection{「正常人被极大高估了」}

原视频在一期叫「ADHD 算是病吗」的视频里，讲了一个颠覆性的观点（BV1PLXQBUEns）：\term{和 ADHD 对比的时候，正常人往往被极大高估了}。

家长最怕的场景是：「我孩子一个晚上一个单词都搞不定，隔壁小孩一天一百个，这差距不是越来越大？」但原视频说：\term{这个线性积累模型本身就是错的}。知识不是线性累积的——按部就班学了多少年、学得一般的人，多的是。十年后，两个孩子很可能都在中游。

这有两层意思，对学历光谱的两端都成立：

\begin{checklist}
  \item \textbf{对学历中下层}：你当年「追不上别人」的焦虑，可能从头到尾建立在错误的模型上。按部就班不是积累的唯一方式，甚至不是大多数人真正掌握知识的方式；
  \item \textbf{对学历上层}：也别太得意。高学历不能证明你的系统健康，只能证明环境替你做了一部分执行。脚手架一撤，谁裸奔谁知道。
\end{checklist}

\subsection{同一颗大脑，两种剧本}

把两章的画像放在同一张表里，差异一目了然：

\begin{table}[h]
\centering\small
\begin{tabular}{>{\raggedright\arraybackslash}p{2.6cm} p{6.1cm} p{6.1cm}}
\toprule
 & \textbf{学历上层} & \textbf{学历中下层} \\
\midrule
上学体验 & 成绩还行，靠聪明和被动注意力混过 & 反复失败，被归因为「懒 / 笨 / 态度问题」 \\
自我概念 & 「我很强，但不能暴露我很乱」 & 「我不行，做不到是我的宿命」 \\
主要燃料 & 死线肾上腺素 + 超专注红利 & 即时刺激 + 兴趣上头 \\
崩塌方式 & 燃尽循环、越成功越空心 & 回避一切、机会滚雪球式流失 \\
自救重点 & 先建系统，停止靠肾上腺素续命 & 先修自我概念，攒「我能做到」的证据 \\
\bottomrule
\end{tabular}
\end{table}

\subsection{中间态：大多数人都在这条光谱上}

当然，现实里绝大多数人不在两个极端，而在中间：学历不低也不高，症状不重也不轻，时而上头时而下沉。这张光谱的意义不是让你对号入座去站队，而是让你\term{认出自己的燃料和崩塌方式属于哪一型，然后对症下药}。

而且两型的经验是可以互相借用的：

\begin{itemize}
  \item 学历上层可以借中下层的「放下」：允许自己蒙混过关、允许自己放弃不重要的事，先混过去再重新产生兴趣；
  \item 学历中下层可以借上层的「系统」：用恰当场、程式性计划这些方法，给自己搭一个没人替你搭的脚手架。
\end{itemize}

\subsection{三个可干预的变量}

最后，把这张光谱的时间轴拉直，会发现决定一个人滑向哪端的关键变量只有三个，而\term{这三个全部可以后天干预}：

\begin{checklist}
  \item \textbf{自我概念}——可以重建。从「我不行」改写成「我这套系统需要不同的运行方式」；
  \item \textbf{环境脚手架}——可以自建。没人替你做执行，就用恰当场、环境布局、外置大脑替你做；
  \item \textbf{早期经验}——可以重写叙事。过去无法改变，但「我当年为什么失败」的解释可以重写：不是我不行，是我一直缺脚手架。
\end{checklist}

\key{剧本不是定死的。同一个算法，过去把你写成了某一型；从现在起，你可以亲手改写这一章。}

\part{自救方法论：顺着大脑参数设计生活}

\section{认知自救：先让自己舒服}

进入方法之前，必须先打地基。因为一个残酷的事实是：\term{方法再好，一边做一边骂自己「我怎么这么懒」，方法也会失效}。

ADHD 自救的顺序，永远是认知在前，操作在后。这一章先把「舒服」这件事说透。

\subsection{从「我有病」到「我有另一套参数」}

你大概在无数个深夜问过自己：我是不是太懒？我是不是没救了？我是不是根本不适合活着这件事？

把这些问题翻译一下，它们都建立在同一个前提上：\term{「我应该像正常人那样运转，而我坏了」}。这个前提一旦松动，整座山都跟着松。

换个说法试试：你不是坏了，你是\term{参数不同}。参数没有对错，只有适配。GPU 跑不了单线程是「缺陷」吗？不是，是它被用错了地方。你过去所有「失败」，都是这套参数在错误的运行环境里的正常结果——不是人品问题，不是能力问题，是\term{不匹配}。

\begin{insight}
病耻感是最贵的内耗。它让你把大量的执行力花在「证明自己没坏」上，而不是花在「把自己用好」上。放下「我有病」三个字，是自救的第一个动作。
\end{insight}

\subsection{停止道德账本}

ADHD 群体有一种特有的内耗：随身携带一本道德账本，每拖延一次，就记上一笔「我又懒了」「我又辜负了期待」。到晚上临睡前，把账本翻出来，一边数一边骂自己。

问题在于：\term{自责不是动力，是阻力}。你花在自责上的精力，够你做三件事了。更糟的是，自责会推高皮质醇——还记得第三章吗？皮质醇一高，启动就锁死，于是你陷入「拖—骂—更拖」的死循环。

\begin{checklist}
  \item 把「我怎么这么懒」改成「今天的燃料又不对路」——前者是审判，后者是观察；
  \item 把「我必须今天完成」改成「我今天先做五分钟」——前者是警报，后者是入口；
  \item 每天晚上，把道德账本换成「今天大脑运行日志」：什么时段状态好、什么刺激管用、什么环境有害——这是数据，不是罪证。
\end{checklist}

\subsection{自我关怀不是放纵}

很多人不敢「放过自己」，是因为怕一放就彻底躺平。这里要分清两件事：

\term{放纵}是「今天算了」，是关闭；\term{自我关怀}是「今天这样也是合理的，明天换个办法」，是重新开始。

差别在哪？放纵之后更恨自己；自我关怀之后更有力气。羞耻会锁死执行功能，接纳会解锁。对 ADHD 来说，自我关怀不是奖赏，是\term{生产力}——它是让大脑敢再试一次的前提。

\subsection{接受「我需要外部脚手架」}

眼镜不可耻，外骨骼也不可耻。

别人靠内建自控（前额叶直接指挥），你靠外置结构（环境、清单、提示、陪伴）。只要结果一样——作业交了、事情做了、生活没塌——\term{谁在乎结构装在大脑里还是装在环境里}？

很多学历上层的 ADHD 在这里卡住：太要强，觉得「靠环境布置来让自己做事」是作弊，显得自己不行。恰恰相反，\term{懂得布置环境，才是对这套算法最深的尊重}。把最重要的文件放桌面最显眼处、把手机放另一个房间、让记录一直开着——这些不是逃避，是工程。

\subsection{允许「蒙混过关」和「放弃」是合法选项}

这是原视频里最有争议、也最被低估的一招（BV1PGcZzsE9K，BV1Y2wizyEua）：对 ADHD 来说，\term{「蒙混过关」和「放弃」应该是合法选项}。

先说放弃：有些事你实在拖太久了，而且它对生活没有绝对影响，那就干脆决定不做。原视频的观察是：\term{你决定不做了，之后重新产生冲动的概率反而更高}——因为「必须做」的警报撤了，大脑不再对抗，兴趣才有机会重新长出来。先混过去，再放下，再重新自发产生兴趣，这个循环比被任务死死压住健康得多。

再说蒙混过关：接到一个外界派下来的线性任务，别从「一丝不苟地做下来」开始（那对 ADHD 是死刑），先想「怎么把它蒙混过关」——这个过程反而会产出很多想法，再迭代成能交的版本。注意，这招涉及社交羞耻（「我一开始就在想蒙混过关，会不会显得我很坏」），心理负担极大，没几个人愿意用，但它确实是存在的、有效的选项。

\begin{notehot}
这两招的真正用途，不是让你变油滑，而是\term{给你腾出重新启动的空间}。对一个启动困难的大脑，「允许不做」和「允许烂做」是两把钥匙：钥匙插进去，门才开得动。等重新有了兴趣和能量，再好好做，那时候才是真质量。
\end{notehot}

\subsection{一个微妙的平衡}

当然也要说清楚：以上都不是「给自己的借口」——借口是逃避之后的自我安慰，这些是\term{给启动让路的工程手段}。区别只有一个检验标准：\term{用完之后，你是更无力了，还是更有力气了}。更无力，那是借口；更有力气，那是方法。请用这个标准，而不是道德标准，来判断自己。

\key{舒服不是躺平，舒服是停止和重力对抗。认知上彻底放下「我应该像正常人一样」，方法才能落地。}

\section{操作自救：一整套 ADHD 原生方法库}

认知打完了地基，这一章是整份讲义的操作核心。这里的每一个方法，都不是「通用自律技巧」的换皮，而是\term{顺着 ADHD 算法本身设计出来的}——它们不要求你提高意志力，而是要求你降低意志的需求。

先给一个大原则：\term{顺着算法，不逆着算法}。所有方法都可以归结为三句话——

\begin{checklist}
  \item \textbf{以广促深}：广度算法便宜，那就让它便宜着干活，别逼它走深度；
  \item \textbf{以被动代主动}：主动切换贵，被动切换便宜，那就多制造被动，少依赖主动；
  \item \textbf{以环境代意志}：意志耗能贵，环境布局便宜，那就把意志活干的事交给环境。
\end{checklist}

\subsection{方法一：恰当场 / 抽象轨道}

这是整套方法的底座（BV11ucZzxEkx）。核心动作只有两个：

\term{第一，把一堆事情，在内心合并成「一件事」。}你这学期有十本书要读，那这十本书就是一件事：「我这学期该读的东西」。A 项目、B 汇报、C 考试，凡是你觉得「同时符合当下和长远」的，都并进一个\term{恰当场}，然后告诉自己：\term{我今天只做一件事}。

\term{第二，内部随便跳，但不跳出轨道。}念头从 A 跳到 B 再跳回 C，无所谓——外部看你在频繁切换，内部看你一直在这条宽轨道上运行。这避开了最贵的东西：主动任务切换（那意味着抑制当前冲动）。你跳的是被动的、便宜的、自发的念头。

范围怎么定？\term{一开始要定得大一些}，大到「下至小学数学、上至前沿研究」都装得下的程度。不用担心因此整天做简单事——ADHD 大脑本身追求新奇，简单的东西无聊得更快，它自己会往深处走。

\begin{insight}
一句话：\term{把「必须做 ABCD」的排序，换成「我只在这一片场里随便蹦」}。排序是高耗能的抑制过程，蹦是低耗能的自发过程。让大脑觉得「我是自由的」，但自由的范围恰好是你定的。
\end{insight}

\subsection{方法二：程式性计划：计划做 A，不限定内容}

ADHD 完全可以做计划，但要做\term{程式性计划}（BV16VckzqEMp）：计划「下午做 A」，但不限定任何具体内容。

于是下午就变成：随便翻 A 的资料、书、视频，翻到哪个有兴致就停一下。\term{只要翻得足够多，激活兴趣的概率总是存在的}——哪怕一下午没写出什么「正事」，你肯定翻到了某个让你眼睛一亮、顺手写了两笔的点。

这里有个心理坎必须过：你怕这样太散、不系统，怕「正常人一下午按部就班看了很多，而我一直在乱看，差距越来越大」。原视频的回答非常直接：\term{不用担心，我在清华待了这么多年，知道正常人一下午也学不了多少}。线性积累模型的错误，我们第六章已经拆过了。

注意一个要点：\term{A 的范围要足够大}，大到你可以随便跳。计划「下午学编程」而不是「下午学第一章」，这样跳到数学、跳到工具、跳到相关想法都不算出轨。

\subsection{方法三：毛点计划：计划只是锚点}

ADHD 做计划，最常见的失败方式是：把计划当执行清单，做不完就自责。原视频给了一个更聪明的定位（BV1Jj8vzWEAQ）：\term{计划只是头脑里的毛点（锚点），不是执行路径}。

你有一个计划，就像心里种了一个钩子。之后在日常里，你会不断布置环境、制造刺激，让一些「啪的一下」的想法冒出来——\term{一拍脑袋冒出来的想法，往往正好和你的计划有关}。这时候立刻去行动。

为什么「啪的一下」很重要？因为\term{这个想法冒出来的时候，你的多巴胺是够的}。启动困难的大脑，最缺的就是这一下。原视频甚至回忆自己早年太严肃地对待这类想法——总想着「这想法靠不靠谱、别人做过没有」，结果把最有动力的启动机会全否定了，退回枯燥的读文献轨道，越想越不想动，陷入恶性循环。

\begin{checklist}
  \item 计划 = 毛点：想得足够远，能激活就算成功，不要求绝对可控；
  \item 环境 = 触发器：文件分类、桌面布局、收藏夹，都朝着「让相关想法冒出来」布置；
  \item 行动 = 追着「啪的一下」走：冒出来就做，别自我否定。好的想法，常常不是在按部就班里迭代出来的，是在兴奋的尝试里突然冒出来的。
\end{checklist}

\subsection{方法四：框架欺骗：撤掉「必须」警报}

这是对付「排异反应」的经典招数（BV1WBPezFEqE）。

你可能有这种体验：心里一旦确定「接下来两小时我要做 A」，身体立刻产生抵抗——想睡、想刷手机、想逃跑。原视频的解释是：\term{「必须做 A」这句话，在神经层面是一个警报}——系统预判「即将进入高耗能模式」，于是提前启动防御，表现为抗拒、拖延、犯困。

解法是\term{框架欺骗}：告诉自己「接下来两小时干啥都可以」，但实际行为局限在 A 相关范围内——你随便翻 A 的资料，只是心里没有「必须」。警报撤了，防御关了，大脑进入开放状态；一旦某个点恰好点燃了兴趣，系统就会\term{带着兴趣自动倒向 A}。

\begin{insight}
这个策略的精妙之处：\term{你给了大脑自由的幻觉，却没有偏离目标}。自由状态下的大脑更容易被激活，而激活的那个点，正好在 A 的范围里。用自由骗过防御，用范围保住目标。
\end{insight}

\subsection{方法五：布局神经元：环境布局代替意志力}

这是原视频的另一个核心概念（BV1MWwJzxEh5，BV13HoGYiEgF）：\term{主动注意力你控制不了，但你可以布置环境，改变你的「猎物偏好」}。

具体操作：

\begin{checklist}
  \item 最重要的一篇文章 / 文件，单独占一个屏幕，占满，别让一堆窗口跟你抢视线；
  \item 正在写的记录一直开着，随时可以落笔；
  \item 关键资料放在「走神后三秒内能切到」的位置，而不是藏在文件夹深处；
  \item 反过来，容易带偏你的东西，从视野里拿走——降低它们的暴露频率，就等于降低它们在你注意力竞争里的权重。
\end{checklist}

原理：你没法命令自己「想到正事」，但可以让「正事」在你的视野里反复出现，于是念头在竞争时更可能落到它上面。这就是\term{以环境代意志}——你改不了注意力，但你能改环境给注意力喂什么。

\subsection{方法六：冲动引导成网，随手记录}

第三章讲过冲动理性。这里给它的操作化（BV1apQ2ByEN3）：

\term{冲动来了，别硬压，也别全放。}硬压，抑制耗能爆表；全放，方向完全失控。正确姿势是：\term{记录它，然后继续}。冒出一个想法，用 Obsidian 随手记两行（甚至只是几个自然语言的关键词），清掉一部分工作记忆，然后跳回原来的事。这个想法是种子，之后它可能自己长成深刻的问题。

关键细节：\term{首次冲动大多没啥用，但你别懒得记}。因为二次冲动、三次冲动往往是从首次冲动的残留上长出来的，越记越长、越记越深。原视频原话：\term{「记忆不得想法，就是这些首次冲动上长出来的二次冲动」}——很多好问题，就是靠「随手记一下」喂出来的。

方向怎么控？不是靠压，是靠\term{调节环境暴露概率}：你同时想提高 A 的比例、降低 B 的比例，那就减少环境里引发 B 想法的刺激，多布置引发 A 想法的刺激。冲动已经来了，就释放掉；没冲动的时候，就做 A，并布置更利于 A 自然激发的环境。

\subsection{方法七：隐形回血：管理分心，而非消除分心}

第三章讲过隐形回血。这里补操作细节（BV16hxrzYEpT）：

分心分两类，区别对待：

\begin{checklist}
  \item \textbf{无新内容的冲动}（突然想刷手机、开无关网页）：用身份认同压下去——「我是一个气度沉稳、能按顺序行事的人」，这一句就有用；
  \item \textbf{有价值的创新性想法}（冒出一个新的连接、一个猜想）：立刻记下来或问 AI——这是回血，是好想法的种子，也是启动门槛的降低剂。
\end{checklist}

为什么这算回血？因为全神贯注的效率，会被事后累抵消——硬撑两小时敲完代码，之后躺三小时；允许有价值的分心，四个小时敲完，中间产出一堆想法，还不太累。\term{后者是更划算的}。别当电动车的司机，却告诉车「你要跑得更快，所以不准充电」。

\subsection{方法八：唤醒度三档管理 + 假性疲劳处理}

第二章讲过唤醒度三档。这里把它变成一个可执行的日常流程（BV16mXnBxEE1，BV1MWwJzxEh5）：

\begin{checklist}
  \item \textbf{高唤醒}：做线性任务、写重要的东西，趁热打铁；
  \item \textbf{中唤醒}：随手翻找——电脑里、书里、Obsidian 里，看看有什么能激发兴趣；或做点本身不需要高唤醒的事；
  \item \textbf{低唤醒}：躺下，允许大脑自己放想法。用手机里「方便思考且你觉得有意义」的内容，或者直接问 AI 几个问题，重新唤醒。别硬撑，也别内疚——这不算偷懒，这是这套系统的重启方式。
\end{checklist}

处理假性疲劳的口诀是：\term{想法上高自由度，操作上低自由度}。

\begin{checklist}
  \item \textbf{想法上高自由度}：告诉自己「今天不一定要写这个，随便偏出去写点别的，也是成果」——撤掉「必须」警报；
  \item \textbf{操作上低自由度}：要写的文件放桌面最显眼处，记录一直开着——让任何想法都能在几秒内转向它。
\end{checklist}

\subsection{方法九：音乐策略}

「ADHD 学习能不能听音乐」——原视频给了比「能/不能」更细的答案（BV1T1XSBnE6w）：

\begin{checklist}
  \item \textbf{机械熟练任务}（整理、搬运、抄录、打码）：放音乐。不放音乐时用来防分心的自我监控成本，可能比放音乐还高；音乐（尤其是纯音乐）还能占住想跳出的念头，收窄跳跃范围；
  \item \textbf{随意翻找内容}：放不放影响不大；
  \item \textbf{看似机械、但需要预测每一步的任务}（不熟的步骤）：关音乐。因为你的大脑每一步都在预测，那需要主动注意力，别和音乐抢带宽。
\end{checklist}

核心逻辑：机械任务上争分夺秒没意义——它们永远只是准备工作，关键在「做」，不在「快」。把省下来的意志留给真正重要的任务。

\subsection{方法十：模糊分类：给广度搜索装一个范围}

ADHD 的搜索是广度优先的（BV1xpLFzaECh）。当头脑里缺乏「粗糙的分类」，搜寻范围过大，大脑就会「一团浆糊」——看理工科内容尤其容易「脑内温暝、像烧开水一样」。解法是\term{模糊分类}：

\begin{checklist}
  \item 给一本书的十三四章，自己概括成四五个大类，按你的方式起名字（「攻防法服」那种级别的粗糙分类）；
  \item 再往上，给整个学科建大目录——根据你读的文章和脑里的想法，把学科分块；
  \item 分类的名字必须概括特征，不能瞎起；但边界模糊没关系，模糊正是它的优点。
\end{checklist}

为什么模糊比严格好？\term{严格化要你平白多记很多东西，还破坏了强弱顺序}。而粗糙分类能一眼看出「这两个东西同属一类，所以能用同样的方法处理」——省下的算力是巨大的。古人把东西分成阴阳五行、地水火风，对古人来说可能真的有用，因为它限定了搜寻范围。

\subsection{方法十一：慕容水战法（或战或亡）}

这是原视频最戏剧性的一招（BV1zDpkzZEtX）。先说结论：\term{慢性压力是 ADHD 的天敌，急性压力反而释放认知资源}。

慢性压力（做不完要挨骂、要失业、做不出来）推高皮质醇，皮质醇一高，启动锁死——你陷入「这任务做完也没啥用 + 做不完要被评判」的死循环，五点就想睡。而\term{无路可退的急性压力}，是另一套机制：当大脑只剩下「或战或亡」一个选项，所有「做不做、什么时候做、怎么做更好」的内耗瞬间消失，目标变得无比清晰——就是眼前这件事。

原视频把它叫\term{慕容水战法}：想象自己像慕容垂那样，七十岁、久病缠身，但这一次出征没有退路——「此刻只负责把眼前这一仗打完，不再想后续」。这个框架一旦进入，认知资源反而被释放，内耗消失，反而能干活了。

\begin{notehot}
必须强调：\term{这招不能常用}。它有点「先伤己、再伤人」的性质，是伤敌一千自损八百的绝地求生。它的价值在于：当你被慢性压力困死的时候，知道还有这么一把钥匙——把「想太多」换成「只打眼前这一仗」，哪怕只是假想地进入这个状态，也能达到部分效果。自我欺骗只要能起实际的好效果，从神经学角度看就是真实。
\end{notehot}

\subsection{方法十二：工具层}

原视频专门推荐过一套工具组合（BV1xMonY7EGD，BV16hxrzYEpT）：

\begin{checklist}
  \item \textbf{Obsidian}：打开文件快、可折叠、有关系图谱——天生适合「想法冒出来随手记、跳回来再找」的工作流；
  \item \textbf{虚拟桌面划分}：桌面一放最重要的一篇文章（单内容占满），桌面二放 Obsidian，桌面三放常用查询网站。主动切过去是自发的、可以的；被 Alt+Tab 弹出一堆窗口被动干扰，才是要避免的；
  \item \textbf{AI 即时反馈}：低唤醒 / 清醒睡眠状态时，你的符号语言处理能力掉线，但自然语言还在。把数学对象、猜想、想法用自然语言讲给 AI，当下就拿到反馈——原视频认为这对 ADHD 极其友好，是过去没有的「即时奖励」。
\end{checklist}

\subsection{方法十三：游戏与刺激的重新管理}

最后是一点生活管理（BV1U1XbBfEtE）。ADHD 经常在游戏上栽跟头，但原视频的建议很反直觉：

\term{不要在一天内规划「游戏两小时」。}这种精确配额带来的结果往往是：玩得不够尽兴、一超标就觉得「反正超标了，干脆玩到吐」，变成一天十几个小时、连续几个月。改成\term{直接滑一两个整天去玩}——把游戏当成一件被允许的事，而不是要偷摸超标的事。意志耗能降下来，反而容易收手。

对睡眠的态度也一样（BV16mXnBxEE1）：不必纠结「必须马上睡着」。晚上到点躺下就行，睡不着时头脑也在休息；躺床上用手机问 AI、看有趣内容，不叫没休息——对 ADHD 来说，那就是休息的一种。

\subsection{一个通用的启动流程}

把以上十二个方法串起来，就是一条可复制的启动流水线：

\begin{enumerate}[leftmargin=2em,itemsep=2pt]
  \item 先处理无聊：机械任务放音乐，中低唤醒就翻找激发，别硬坐；
  \item 用框架欺骗撤掉「必须」警报：告诉自己「接下来两小时干啥都行」，实际只翻 A 相关；
  \item 在恰当场里被动切换：念头冒出来就跟，跳出轨道的不跟；
  \item 有价值的想法：随手记录或问 AI，那是回血和种子；
  \item 假性疲劳：动态调自由度——想法上放开，操作上把文件放最显眼处；
  \item 被慢性压力困死时：换一次慕容水战法，只打眼前这一仗。
\end{enumerate}

\key{这一章没有一个方法是「逼自己」的。它们全是同一句话的变体：别硬改大脑，去改它运行的场地。}

\section{分层策略：从你的位置出发}

方法库是通用的，但你的起跑线不一样。这一章把前面的所有东西，按学历光谱的两端分别组装成两套不同的自救方案。

\subsection{学历上层：从「撑住」到「可持续」}

你的问题不是「动不了」，是「靠肾上腺素动，快烧完了」。所以你的自救关键词只有一个：\term{拆弹}——拆掉那台让你撑到现在、也快把你烧穿了的引擎。

\begin{checklist}
  \item \textbf{停止等死线}。死线是你的拐杖，也是你的癌症。它让你每次都能「爆发」，也让你每次都在「透支」。从今天起，选一件重要但不紧急的事，用程式性计划每天喂它一点，喂够一个月——你会第一次体会到「不是被逼着做完」是什么感觉；
  \item \textbf{建立不依赖紧迫感的系统}。紧迫感是你的多巴胺来源，但它是高利贷。用恰当场 + 环境布局替代它：把项目文件放桌面、记录开着、允许自己每天只在恰当场里蹦跶，让「被动推进」取代「死线冲刺」；
  \item \textbf{允许自己不完美}。完美主义是拖延的遮羞布——「我要一次做完美」翻译过来就是「我害怕交出去被评判」。先交一个烂版本，再迭代（这就是第八章的蒙混过关法在你的场景里的合法用法）；
  \item \textbf{警惕自我药物治疗}。咖啡因、酒精、报复性熬夜、报复性刷手机，都是慢性透支。它们不是坏习惯，是你燃料系统告急的警报。先解决燃料问题（唤醒度管理、睡眠、运动），再谈戒断；
  \item \textbf{把超专注引导到长线项目}。你有超专注的红利，别让它全耗在即时刺激上。用布局神经元，让「啪的一下」更多地落在你的研究、创作、事业上；
  \item \textbf{承认自己需要脚手架}。哪怕别人看不出来，哪怕你觉得「我这种人怎么会需要帮助」。主动找 accountability（一个陪你做事的人、一个定期汇报的对象、一个公开承诺），这不丢人，这是工程。
\end{checklist}

\begin{insight}
学历上层的最大敌人是「面子」——「我看起来一切正常，怎么好意思说我不行」。但请记住：\term{你不是从「正常」滑向「不正常」，你是从「靠透支维持的正常」走向「靠系统维持的正常」}。后者更慢，但不会再崩。
\end{insight}

\subsection{学历中下层：先修自我概念}

你的问题不是「方法不够」，是「信念不够」——你不相信自己能做到，于是再好的方法都会在第一次失败时被你当成「果然不行」的证据。所以你的自救顺序完全不同：\term{先修自我概念，再谈方法}。

\begin{checklist}
  \item \textbf{小赢螺旋}：每天做一件「能完成」的小事，攒「我能做到」的证据。不是「背一百个单词」，是「今天翻开了书，看了三页」，甚至只是「把桌面整理干净」。证据比道理有用一万倍——道理说服不了被反复否定过的自我概念，证据可以；
  \item \textbf{别跟主流赛道死磕}。你不需要在「别人定义好的线性赛道」上证明自己。选择高刺激、即时反馈、行动导向的工作与学习方式——手艺、销售、技术操作、自由职业、创业，都是广度算法的好赛场。学历低不是你的污点，是你提前逃离了不匹配的赛道；
  \item \textbf{利用免费资源}。公开课、免费工具（Obsidian）、AI——现在一个手机就能搭出十年前博士才能有的外置大脑。你的脚手架，可以自己建；
  \item \textbf{承认并利用隐藏技能}。动手能力、口才、直觉、危机处理、看人准——这些学校不给你打分，但市场会给。把「不务正业」重新命名为「技能」；
  \item \textbf{允许自己从「放弃」开始}。对不重要的事，先混过去、再放下、再重新产生兴趣。别让「我必须把这件事做完」的执念，堵死你重新启动的全部可能。
\end{checklist}

\begin{insight}
学历中下层最需要听的一句话：\term{你不是不行，你只是一直没被给过脚手架和第二次机会}。现在，脚手架可以自己搭，第二次机会可以自己给。从一个小赢开始，让「我能做到」的证据滚雪球。
\end{insight}

\subsection{通用的一步：身份重建}

无论你在光谱的哪一端，最终都会走到同一步——\term{身份重建}。

把这几个句子抄下来，贴在每天看得到的地方，它们是你的新身份证：

\begin{checklist}
  \item 我不是懒，我是\term{燃料不匹配}——我的启动器不听「重要」的话，只听兴趣、新奇、挑战、紧迫的话；
  \item 我不是不能专注，我是\term{主动专注太贵}——所以我要学会让被动注意力落在正事上；
  \item 我不是拖延，我是\term{被「必须」锁死了}——撤掉警报，我就能动；
  \item 我不是执行力差，我是\term{一直在用最贵的模式干活}——改用便宜的模式，我比谁都持久；
  \item 我是一个\term{用不同燃料运行的人}，而且我正在学会怎么加油。
\end{checklist}

\key{从「我是懒人」到「我是一个用不同燃料运行的人，而且学会了怎么加油」——这一步迈过去，前面的所有方法才算真正长在你身上。}

\part{更深一层的思考：与这个世界相处}

\section{思考：一个更舒服的看待世界的角度}

方法讲完了。这一章不讲操作，讲一个更舒服的看待世界的角度——把前面那些散落的认知，拼成一张完整的世界观。

\subsection{为什么学校会惩罚 ADHD}

先说一个很多人没想明白的问题：为什么 ADHD 在学校里总是吃亏？

答案不是「学校更难」，而是\term{学校默认深度优先算法}——一节课四十五分钟、一章接一章、线性积累、一次没跟上就永远跟不上。这套运行模式恰好踩在 ADHD 算法最贵的那块成本上。上课听不懂，浪费时间，产生挫败感；搜寻范围过大，缺乏粗糙分类，脑内一团浆糊；突发兴趣转不过去，想法无法深入发展。

\begin{insight}
注意这里的关键：\term{这不是你的错，是环境不匹配}。你在课堂里「失败」，和你拿着 GPU 去跑单线程渲染失败，性质是一样的——工具没错，是任务和工具不匹配。把「我上课听不进去」翻译成「我在一个不适合我的运行环境里」，你就不会再用「我笨」「我没救」来惩罚自己。
\end{insight}

\subsection{「轴的程度」与「适不适合」：维度不匹配不是病}

原视频讲过一个小案例，非常能说明问题（BV13HoGYiEgF）：一个能考上北大光华的学生，高等数学考三四十分，怎么教都不会。你很难说他是笨——能上光华的人，怎么会连高数都学不会？更合理的解释是：\term{某个维度不匹配}。

原视频还有个更狠的类比：人的「轴的程度」各不相同——轴到九十九分的人不适合做数学，轴到四五十分的人也不适合，八十分左右最适合。那是不是该给「高轴」的人判一个「高轴综合症」、给「低轴」的人判一个「低轴综合症」？这当然是搞笑——\term{谁规定你生下来就得做数学了}？如果按「适不适合做数学」来给人分类，多分几个维度，每个人都有一堆病。

所以对「我是不是有缺陷」这个问题，一个更舒服的回答是：\term{缺陷与否，取决于它在你生活里造成了多大困扰}。特征 + 损害 = 诊断；没有损害，就只有特征。方法解决得了的困扰，就不是「病」，是「待适配」。

\subsection{疾病观的延伸：方法优先，药物兜底}

顺着这个逻辑，原视频对「ADHD 算不算病」的态度很明确（BV13HoGYiEgF）：如果你靠布局神经元、连接整理那一套手法，依然觉得费劲、难以正常工作，那你确实需要药物，不必硬扛；但如果方法能解决，那就没必要把它当病治。

这里的核心立场是\term{实用主义}：不问「我算不算病」，只问「什么让我运转得最好」。方法能解决的用方法，方法解决不了的，药物和正规医疗是合理且正当的兜底——\term{用工具解决问题，不丢人}。这两条路不冲突，更不是道德高下之分。

\subsection{以广促深：广度算法在长周期里能赢}

ADHD 的广度算法，在短周期、单任务的视角下是「缺陷」；但放到长周期、多任务的视角下，它有独特的优势。

\term{深度运算虽然贵，但「以广促深」是成立的}：你了解并领悟很多不同的视角，切换着看问题，最后能触达的深度，不一定比线性深挖的人浅。很多重要发现恰恰来自「把 A 领域的方法搬进 B 领域」——而 ADHD 的广度搜索，天生擅长制造这种跨界连接。

原视频甚至有个大胆的猜想：\term{欧拉可能是 ADHD}。欧拉晚年失明了，依然能不停口述写作、产出海量论文——因为他的写作是「默认模式」，靠被动注意力驱动，根本不需要控制。\term{这是广度算法被逼到极限之后的样子：不需要开关，一直在产出}。

\begin{insight}
所以「ADHD 是不是缺陷」，最后的答案是这样的：它是一套在当下环境里处处吃亏、但在长周期里藏着独特优势的算法。前期的关卡难打，后期的优势真实存在。你不需要为「广度」道歉，你需要的是把它用对地方。
\end{insight}

\subsection{关于「正常人」：他们也被极大高估了}

整份讲义里，我们反复提到「普通人」「正常人」。最后要提醒一句：\term{他们也被极大高估了}。

你羡慕的那个人，可能只是恰好活在适合他的环境里；你觉得自己「连普通人都不如」，可能只是你拿自己最弱的运行模式，去比人家最强的运行模式。多分几个维度，每个人都有「病」，也每个人都有「天赋」。

原视频用一句很朴素的话收过这个主题：\term{不要老想着别人能做到，你也一定能做到}——这句话反过来也成立：\term{别人做到了，也不代表你必须用同一条路做到}。你可以最后做出好的作品，但过程不需要和别人一样。孙子那句话用在这里很合适：\term{能屈能伸，方为丈夫}。

\subsection{给 ADHD 群体的一句话}

把这一切收束成一句给同路人的话，原视频的说法是：

\term{先混过去，再放下，再重新自发地产生兴趣。相信顺应冲动，也能做出很多事情。}

你不是需要被治愈，你是需要被读懂。这份讲义，就是读懂自己的开始。

\section{终章：一页纸的每日自救清单}

前面十章，浓缩成一张可以贴在墙上的纸。

\subsection{每天的三件事}

\begin{checklist}
  \item \textbf{晨间}：定下今天的「一件事」（恰当场）。把最重要的文件放到最显眼处，记录打开。别强求早起，晚上到点躺下就行；
  \item \textbf{午间}：低唤醒就翻找激发或躺下等想法冒出来，别硬坐；机械任务放音乐；冒出来的有价值想法，随手记或问 AI；
  \item \textbf{晚间}：允许自己在恰当场里被动切换；用「接下来干啥都行」骗过「必须」警报；不给自己定「今天只能玩两小时」这种精确配额。
\end{checklist}

\subsection{每周的一小时}

\begin{checklist}
  \item 给正在学的一个学科或一本书，建一次\term{模糊分类}——四五个大类，各起个名；
  \item 做一次\term{连接整理}：整理文件、收藏夹，把散落的东西连起来，删掉不满意的旧稿；
  \item 复盘：这周哪些环境、哪些毛点最容易激活你？记进大脑运行日志。
\end{checklist}

\subsection{遇到卡壳时，按这个顺序试}

\begin{enumerate}[leftmargin=2em,itemsep=2pt]
  \item 我是不是在\term{无聊}？——先处理无聊：放音乐、翻找激发、换个相关的东西看；
  \item 我是不是被\term{「必须」锁死}了？——换成「接下来两小时干啥都行」，实际只碰 A 相关；
  \item 我是不是在\term{低唤醒}状态？——别硬撑，躺下等想法冒出来，或用手机 / AI 唤醒；
  \item 我是不是\term{假性疲劳}？——撤掉「必须做 X」的念头，把文件放最显眼处，随便写几笔；
  \item 我是不是被\term{慢性压力}困死了？——换一次慕容水战法：此刻只打眼前这一仗。
\end{enumerate}

\subsection{永远记住的四句话}

\begin{checklist}
  \item 你不是不能专注，你是\term{主动专注太贵}——让被动注意力落在正事上；
  \item 你不是懒，你是\term{燃料不匹配}——你的启动器听兴趣、新奇、挑战、紧迫的话；
  \item 学历不是能力，是\term{脚手架}——缺了可以自己搭，别拿它审判自己；
  \item 自救不是变成别人，是\term{搞懂自己，然后顺着自己活}。
\end{checklist}

\subsection{最后的话}

这份讲义读到这里，你大概已经知道自己是哪一种了：可能是那个靠聪明和死线撑了二十年的「高功能者」，可能是那个被「我不行」困了很多年的「未被发现者」，更可能是夹在中间的、大多数。

无论哪种，终点都一样：\term{你不是坏掉的，你是不同频率的}。

从此你不再问「我为什么做不到」，而是问「我这套系统今天吃什么燃料、走哪条轨道、在哪个场地里跑得最好」。前者是审判，后者是工程。把审判换成工程，这份讲义就没有白读。

去把今天的那件「一件事」定下来吧。就从现在开始。

\appendix

\section*{附录 A：给「怀疑自己是 ADHD」的自检线索}

这份清单不是诊断工具，它只做两件事：\term{帮你发现自己}，以及\term{帮你组织语言}（如果你打算去和专业人士聊）。逐条对照，不必每条都中——前面说过，伴随特征个体差异极大。

\subsection*{启动与任务}

\begin{checklist}
  \item 启动一件「没有兴趣、但有回报」的事，非常困难，常常拖到死线；
  \item 一旦没有外界任务压着，反而能自由地做自己的事、发展自己的想法；
  \item 对「必须做 X」这类命令有强烈的排异反应：犯困、想逃、刷手机；
  \item 任务「做完了没有」「做到哪一步了」很难自己追踪，需要别人提醒。
\end{checklist}

\subsection*{注意力}

\begin{checklist}
  \item 能连续数小时沉浸（超专注）在自己感兴趣的事上，但没法「按需专注」；
  \item 上课、开会、读长文时极容易走神，一走就拉不回来；
  \item 主动切换任务很痛苦，被动冒出念头跟着跑反而舒服；
  \item 经常「人醒着，但脑子像睡着了一样」：低唤醒、假性疲劳、午休后更累。
\end{checklist}

\subsection*{时间与组织}

\begin{checklist}
  \item 经常忘时间、错过约定、低估任务耗时；
  \item 生活层面功能崩塌：堆积的快递、没回的微信、乱到失控的房间；
  \item 需要靠死线、紧迫感才能启动；越没有压力越动不了。
\end{checklist}

\subsection*{情绪与冲动}

\begin{checklist}
  \item 对拒绝、批评、被评判高度敏感，容易过度反应或长期内耗；
  \item 容易上头也容易下头：三分钟热度、频繁开始又放弃；
  \item 用咖啡因、酒精、游戏、短视频、暴食等即时刺激自我调节；
  \item 心里有一本道德账本：长期用「我怎么这么懒」审判自己。
\end{checklist}

\subsection*{怎么使用这份清单}

\begin{enumerate}[leftmargin=2em,itemsep=2pt]
  \item 划掉那些和你无关的条目——这是正常的，伴随特征本来就因人而异；
  \item 找出 3--5 条对你生活\term{造成实际损害}的条目——那才是你的主线任务；
  \item 如果这些损害已经严重影响工作、关系或生活，去正规医院精神科或临床心理科聊聊。这不丢人，也和继续使用这份讲义不冲突；
  \item 无论是否确诊，第二、三部分的认知与方法都适用——它们依赖的是「你的大脑运行日志」，不是一张诊断书。
\end{enumerate}

\section*{附录 B：讲义黑话速查表}

\begin{center}
\small
\begin{tabular}{>{\raggedright\arraybackslash}p{3.2cm} p{12.0cm}}
\toprule
\textbf{黑话} & \textbf{含义} \\
\midrule
算法观 & 把 ADHD 理解为「广度便宜、深度昂贵；被动便宜、主动昂贵」的算法，而不是一处故障 \\
椭圆 & 脑内无意识念头的代称；超专注 = 某个念头占满注意力空间 \\
超专注 & 沉浸数小时、外界全屏蔽的状态；电源开关不在自己手上，但可引导 \\
归并效应 & 标签被广泛使用后，非必然的伴随特征也被贴上来的现象 \\
伴随特征 & 除执行功能缺陷外的附加特征，个体差异极大，甚至互相矛盾 \\
执行功能 & 维持线性集中 + 抑制其他念头（意志）；两者是一体的 \\
意志耗能 & 压住一个念头去干别的，可能比做那件事本身更费电 \\
恰当场 / 抽象轨道 & 把一堆事在内心合并成「一件事」，每天只在这一条宽轨道里被动切换 \\
程式性计划 & 计划「做 A」但不限定内容，随便翻，翻多了总会激活兴趣 \\
毛点计划 & 计划只是头脑里的锚点，靠环境刺激让「啪的一下」冒出来再行动 \\
框架欺骗 & 告诉自己「两小时干啥都行」，实际限制在 A 范围内，撤掉「必须」警报 \\
布局神经元 & 布置环境改变「猎物偏好」，降低启动门槛，以环境代意志 \\
冲动理性 & 冲动伴随新想法、神经生长与低耗能；顺着冲动调控成「网」而非「线」 \\
时间理性 & 默认「干 A 还是 B 只差切换成本」的普通理性模型，ADHD 上站不住脚 \\
动机悖论 & 社会规范压力推高皮质醇锁死启动；「人类要完蛋」这类念头反而消除压力、释放执行力 \\
隐形回血 & 有价值的创新性分心是回血；管理分心而非消除分心 \\
唤醒度三档 & 高→线性任务；中→翻找激发；低→躺下等想法冒出来再唤醒 \\
假性疲劳 & 缺的是定向能量而非追随性能量；「必须做 X」的念头本身在前额叶耗电 \\
动态自由度 & 想法上高自由度 + 操作上低自由度 \\
模糊分类 & 给学科建 4--5 个粗糙大类以限定广度搜索范围、节约算力 \\
慕容水战法 & 把任务当「或战或亡」的生死战，急性压力反而释放认知资源；不可常用 \\
蒙混过关法 & 外界线性任务先想「怎么蒙混过关」再迭代；心理负担大，少有人愿意用 \\
以广促深 & 广度算法在长周期里能解决深度问题，是贯穿全讲义的总心法 \\
\bottomrule
\end{tabular}
\end{center}

\clearpage
\section*{附录 C：出处对照}

本讲义的思想内核提炼自 B 站 UP 主「阿卡迪萨」（ADHD 野路子）的 24 期视频。以下为主要概念与原视频的对照，方便你回看原作：

\begin{small}
\begin{longtable}{>{\raggedright\arraybackslash}p{4.2cm} >{\raggedright\arraybackslash}p{11.0cm}}
\toprule
\textbf{视频（标题）} & \textbf{本讲义中的相关章节 / 概念} \\
\midrule
\endfirsthead
\toprule
\textbf{视频（标题）} & \textbf{本讲义中的相关章节 / 概念（续）} \\
\midrule
\endhead
\bottomrule
\endfoot
ADHD 是一套不同的算法 & 第一章算法观、椭圆竞争、以广促深；第十章 \\
ADHD 和 ASD 标签的「归并效应」 & 第一章核心特征与伴随特征、确诊 = 特征 + 损害 \\
ADHD 算是病吗？ & 第六章「正常人被极大高估」；第十章疾病观 \\
ADHD 学习的时候该听音乐吗？ & 第八章方法九（音乐策略） \\
ADHD 的执行功能问题及如何调控 & 第二章意志耗能；第八章恰当场、游戏管理 \\
ADHD 醒着但头脑已经「睡着」了 & 第二章低唤醒、唤醒度三档；第八章 \\
为什么午休后感觉更累？ & 第二章低唤醒与被动休息无效 \\
「时间理性」和「冲动理性」 & 第三章；第八章冲动引导成网 \\
游戏中无敌后的无聊现象 & 第三章兴趣引导；「激发场」 \\
「假性疲劳」与「动态自由度调节」 & 第二章；第八章方法八 \\
「程式性计划」但不限定内容 & 第八章方法二 \\
一个有效但没人愿意学的方法 & 第七章蒙混过关法；第九章 \\
应对主动任务切换 & 第二章主动 / 被动切换；第八章恰当场 \\
对「分心」的适当利用 & 第二章无聊优先；分心管理 \\
ADHD 的「动机悖论」 & 第三章动机悖论、一元论视角 \\
认知过载的异质性 & 第八章方法四（框架欺骗） \\
抑制可能比做事更耗能 & 第二章意志耗能；第八章 \\
隐形回血机制 & 第三章；第八章方法七 \\
古代人与现代人的压力；或战或亡 & 第八章方法十一（慕容水战法） \\
抗慢性压力的能力偏弱 & 第二章；第八章方法十一 \\
ADHD 可以做计划吗？ & 第八章方法三（毛点计划） \\
算法和抑制系统的补充（模糊分类） & 第八章方法十 \\
对 ADHD 好用的电脑软件 & 第八章方法十二（工具层） \\
在 B 站搞笑后发现自己好像是 ADHD & 算法观总述；第十章轴的程度、欧拉猜想 \\
\end{longtable}
\end{small}

\vspace{1em}

最后再重申一次，也是原视频一直坚持的态度：\term{以上大多是猜测，个体差异极大}。请把这份讲义当作一张地图，而不是一份处方——地图的作用是让你看清地形，路终究要你自己走，且只有你自己能走。

祝你，从读懂自己开始，慢慢活得顺一点。


\end{document}
